\documentclass[11pt,a4paper]{article}

% ---------- Polices et langue (compilation : XeLaTeX) ----------
\usepackage{fontspec}
\setmainfont{Latin Modern Roman}
\setsansfont{Carlito}
\usepackage{amsmath,amssymb}
\usepackage{microtype}
\frenchspacing
\renewcommand{\contentsname}{Table des matières}
\renewcommand{\figurename}{Figure}
\renewcommand{\tablename}{Tableau}
\renewcommand{\listfigurename}{Table des figures}

% ---------- Mise en page ----------
\usepackage[a4paper,left=2.8cm,right=2.8cm,top=3.1cm,bottom=3cm,headheight=15pt,headsep=0.7cm,footskip=1.1cm]{geometry}
\usepackage{setspace}
\setstretch{1.12}
\usepackage{xcolor}
\definecolor{vert}{HTML}{1F3A2E}
\definecolor{ble}{HTML}{D9A441}
\definecolor{terre}{HTML}{B5532E}
\definecolor{gris}{HTML}{5B6F66}
\usepackage{graphicx,float,booktabs,array,caption,enumitem}
\captionsetup{font=small,labelfont={bf,color=vert},labelsep=period}
\setlist{itemsep=2pt,topsep=4pt}
\usepackage[colorlinks=true,linkcolor=vert,urlcolor=terre,citecolor=vert]{hyperref}
\usepackage{titlesec}
\titleformat{\section}{\Large\bfseries\color{vert}}{\thesection.}{0.6em}{}
\titleformat{\subsection}{\large\bfseries\color{vert}}{\thesubsection}{0.6em}{}
\titlespacing*{\section}{0pt}{1.6em}{0.8em}
\titlespacing*{\subsection}{0pt}{1.2em}{0.5em}

% ---------- Cadre sur toutes les pages (page de garde incluse) ----------
\usepackage{eso-pic}
\usepackage{tikz}
\AddToShipoutPictureBG{%
  \begin{tikzpicture}[remember picture,overlay]
    \draw[line width=1.6pt,color=vert]
      ([shift={(1.3cm,1.3cm)}]current page.south west) rectangle
      ([shift={(-1.3cm,-1.3cm)}]current page.north east);
    \draw[line width=0.5pt,color=ble]
      ([shift={(1.55cm,1.55cm)}]current page.south west) rectangle
      ([shift={(-1.55cm,-1.55cm)}]current page.north east);
  \end{tikzpicture}%
}

% ---------- En-têtes et pieds de page ----------
\usepackage{fancyhdr}
\pagestyle{fancy}
\fancyhf{}
\fancyhead[L]{\small\color{gris}AHNNAOU Youness · CI02}
\fancyhead[R]{\small\color{gris}ENCG Settat · 2026--2027}
\fancyfoot[C]{\small\color{gris}\thepage}
\renewcommand{\headrulewidth}{0.3pt}
\renewcommand{\footrulewidth}{0pt}

\begin{document}

% ======================= PAGE DE GARDE =======================
\begin{titlepage}
\thispagestyle{empty}
\centering
\vspace*{0.4cm}
{\large\bfseries\color{vert} UNIVERSITÉ HASSAN 1\textsuperscript{er} (UHP1)}\\[0.25cm]
{\large\bfseries\color{vert} ÉCOLE NATIONALE DE COMMERCE ET DE GESTION}\\[0.15cm]
{\large\bfseries\color{vert} ENCG SETTAT}\\[0.5cm]
{\color{gris} Filière : \textbf{COMMERCE INTERNATIONAL}}\\[0.2cm]
{\color{gris} Année universitaire 2026--2027}

\vspace{2.2cm}
{\color{terre}\bfseries\large Projet CI02 --- Régression et prévision}\\[0.6cm]
{\Huge\bfseries\color{vert} Prévision de l'indice des quantités de blé importées au Maroc}\\[0.7cm]
{\large Données FAOSTAT (Trade Indices) --- référence naïve, régression Ridge et Random Forest}

\vspace{2.4cm}
\begin{tabular}{@{}r@{\hspace{0.8em}}l@{}}
\textbf{Réalisé par :} & AHNNAOU YOUNESS \\[0.2cm]
\textbf{Numéro Apogée :} & 23009126 \\[0.2cm]
\textbf{Filière :} & Commerce International \\[0.2cm]
\textbf{Professeur :} & Abderrahim Larhlimi \\[0.2cm]
\textbf{Établissement :} & ENCG Settat, UHP1 \\[0.2cm]
\textbf{Année universitaire :} & 2026--2027 \\
\end{tabular}

\vfill
{\small\color{gris} Notebook, données et figures : voir les fichiers du dépôt du projet.}
\vspace{0.3cm}
\end{titlepage}

\tableofcontents
\clearpage

% ======================= 1. INTRODUCTION =======================
\section{Introduction et problématique}

Le Maroc dépend largement des marchés extérieurs pour son approvisionnement en blé. Anticiper l'évolution des quantités importées présente un intérêt direct pour la sécurité alimentaire, la gestion des réserves de devises et l'organisation logistique des ports. Pour un étudiant de la filière Commerce International, c'est aussi un cas concret d'utilisation de la donnée commerciale pour appuyer une décision.

La question de recherche est la suivante :

\begin{quote}
\emph{Comment prévoir l'évolution des quantités annuelles de blé importées par le Maroc à partir de l'historique des échanges commerciaux ?}
\end{quote}

Le périmètre est volontairement restreint afin de rester lisible et vérifiable : un seul pays (le Maroc), un seul produit (le blé), une seule variable à prévoir et une fréquence annuelle. La tâche est une \textbf{régression sur série temporelle univariée}. Un modèle est jugé utile s'il fait mieux qu'une \textbf{référence naïve} (la valeur de l'année précédente), évaluée avec la MAE et la RMSE.

\paragraph{Remarque importante sur la cible.} La variable disponible dans la source n'est pas un tonnage mais un \textbf{indice des quantités importées}, de base 100 sur la période 2014--2016. Toutes les valeurs et toutes les erreurs de ce rapport sont donc exprimées en \emph{points d'indice}.

% ======================= 2. DONNÉES =======================
\section{Données}

\subsection{Source et filtrage}
Les données proviennent de FAOSTAT (domaine \emph{Trade Indices}), à partir d'un export CSV du 19/11/2026 de 18\,240 lignes et 14 colonnes. L'extraction conserve exactement :
\begin{itemize}
  \item zone : \texttt{Morocco} ;
  \item produit : \texttt{Wheat} ;
  \item élément : \texttt{Import Quantity Index (2014-2016 = 100)}.
\end{itemize}
Le filtre retourne 64 observations annuelles, de 1961 à 2024, sans valeur manquante et sans année dupliquée (deux vérifications automatiques dans le notebook).

\subsection{Statistiques descriptives}
\begin{table}[H]
\centering
\caption{Statistiques descriptives de l'indice (64 observations, 1961--2024).}
\begin{tabular}{lrrrrrr}
\toprule
Moyenne & Écart-type & Min. & 25\,\% & Médiane & 75\,\% & Max. \\
\midrule
47,6 & 32,4 & 2 & 25,5 & 43 & 68,25 & 129 \\
\bottomrule
\end{tabular}
\end{table}

L'indice est faible au début de la période (de 2 à 17 sur les huit premières années) et nettement plus élevé sur la fin (de 66 à 129 entre 2015 et 2024). La série est également volatile : par exemple, elle passe de 129 en 2016 à 74 en 2017.

% Pour ajouter les figures d'exploration produites par le notebook, décommentez :
% \begin{figure}[H]\centering
%   \includegraphics[width=0.9\textwidth]{figures/01_evolution.png}
%   \caption{Évolution de l'indice de quantité de blé importée (Maroc).}
% \end{figure}

% ======================= 3. MÉTHODE =======================
\section{Méthodologie}

\subsection{Variables explicatives}
Seules des informations passées sont utilisées : pour l'année $t$, les variables sont les valeurs des trois années précédentes,
\[
y_{t-1},\quad y_{t-2},\quad y_{t-3}.
\]
Les trois premières années sont perdues par construction, ce qui laisse 61 observations de 1964 à 2024.

\subsection{Découpage chronologique}
Le découpage respecte l'ordre du temps (70\,\% / 15\,\% / 15\,\%), sans mélange aléatoire, afin d'éviter toute fuite d'information du futur vers le passé.

\begin{figure}[H]
\centering
\includegraphics[width=0.95\textwidth]{figures/decoupage.png}
\caption{Découpage chronologique entraînement / validation / test.}
\end{figure}

\subsection{Modèles comparés}
\begin{enumerate}
  \item \textbf{Référence naïve} : la prévision de l'année $t$ est $\hat{y}_t = y_{t-1}$.
  \item \textbf{Régression Ridge} : régression linéaire sur les trois retards standardisés, avec pénalisation $L_2$ :
  \[
  \hat{\beta}=\arg\min_{\beta}\ \sum_{t}\Big(y_t-\beta_0-\sum_{k=1}^{3}\beta_k\,y_{t-k}\Big)^2+\alpha\sum_{k=1}^{3}\beta_k^2 .
  \]
  Ici $\beta$ désigne les coefficients \emph{appris}, tandis que $\alpha\ge 0$ est un \emph{hyperparamètre} qui contrôle la force de la pénalisation.
  \item \textbf{Random Forest} : 200 arbres de profondeur maximale 4, graine 42.
\end{enumerate}

\subsection{Choix de l'hyperparamètre sur la validation}
Quatre valeurs de $\alpha$ sont comparées sur la période de validation (2006--2014), le test n'étant jamais utilisé pour ce réglage.

\begin{table}[H]
\centering
\caption{MAE de validation selon $\alpha$ (Ridge).}
\begin{tabular}{rr}
\toprule
$\alpha$ & MAE validation \\
\midrule
\textbf{0,1} & \textbf{19,18} \\
1 & 19,21 \\
10 & 19,64 \\
100 & 23,24 \\
\bottomrule
\end{tabular}
\end{table}

La valeur $\alpha=0{,}1$ est retenue. Les écarts entre 0,1 et 1 sont minimes : le choix est peu sensible dans cette zone, et une pénalisation forte ($\alpha=100$) dégrade nettement l'erreur.

\subsection{Ajustement final et métriques}
Les modèles Ridge et Random Forest sont réentraînés sur l'union entraînement + validation, puis évalués \textbf{une seule fois} sur le test (2015--2024). Pour $n$ observations de test, de valeurs réelles $y_i$ et de prévisions $\hat{y}_i$ :
\[
\mathrm{MAE}=\frac{1}{n}\sum_{i=1}^{n}|y_i-\hat{y}_i|,\qquad
\mathrm{RMSE}=\sqrt{\frac{1}{n}\sum_{i=1}^{n}(y_i-\hat{y}_i)^2},\qquad
R^2=1-\frac{\sum_i (y_i-\hat{y}_i)^2}{\sum_i (y_i-\bar{y})^2}.
\]
La MAE mesure l'erreur typique en points d'indice, la RMSE pénalise davantage les grosses erreurs, et $R^2$ compare le modèle à une prévision constante égale à la moyenne $\bar{y}$ du test.

% ======================= 4. RÉSULTATS =======================
\section{Résultats}

\begin{table}[H]
\centering
\caption{Performances sur le test chronologique (2015--2024).}
\begin{tabular}{lrrr}
\toprule
Modèle & MAE & RMSE & $R^2$ \\
\midrule
\textbf{Ridge} & \textbf{22,47} & \textbf{28,59} & $-0{,}75$ \\
Random Forest & 23,65 & 29,64 & $-0{,}88$ \\
Naïf (année précédente) & 25,90 & 33,21 & $-1{,}36$ \\
\bottomrule
\end{tabular}
\end{table}

Ridge obtient les meilleures valeurs sur les trois métriques. Sa MAE est inférieure d'environ 13\,\% à celle de la référence naïve (22,47 contre 25,90) ; sa RMSE l'est d'environ 14\,\%.

\begin{figure}[H]
\centering
\includegraphics[width=0.62\textwidth]{figures/mae_comparaison.png}
\caption{MAE des trois approches sur le test.}
\end{figure}

\begin{figure}[H]
\centering
\includegraphics[width=0.95\textwidth]{figures/predictions_test.png}
\caption{Valeurs réelles et prévisions sur le test (2015--2024).}
\end{figure}

\begin{table}[H]
\centering
\caption{Prévisions année par année sur le test (points d'indice).}
\begin{tabular}{rrrrrr}
\toprule
Année & Réel & Naïf & Ridge & Random Forest & Erreur abs. Ridge \\
\midrule
2015 & 66 & 106 & 89,25 & 74,21 & 23,25 \\
2016 & 129 & 66 & 69,96 & 73,47 & 59,04 \\
2017 & 74 & 129 & 109,30 & 73,32 & 35,30 \\
2018 & 79 & 74 & 81,13 & 76,08 & 2,13 \\
2019 & 79 & 79 & 106,32 & 72,46 & 27,32 \\
2020 & 113 & 79 & 79,10 & 75,32 & 33,90 \\
2021 & 96 & 113 & 92,47 & 72,97 & 3,53 \\
2022 & 123 & 96 & 92,28 & 82,72 & 30,72 \\
2023 & 120 & 123 & 115,57 & 81,68 & 4,43 \\
2024 & 105 & 120 & 110,11 & 81,68 & 5,11 \\
\bottomrule
\end{tabular}
\end{table}

% ======================= 5. DISCUSSION =======================
\section{Interprétation et analyse des échecs}

\begin{itemize}
  \item \textbf{La référence naïve rate les retournements.} Comme elle recopie l'année précédente, elle se trompe fortement dès que la série change de sens : en 2016, elle prévoit 66 pour une valeur réelle de 129, puis elle prévoit 129 en 2017 pour une valeur réelle de 74.
  \item \textbf{La Random Forest reste trop plate.} Ses prévisions restent entre 72 et 83, alors que les valeurs réelles atteignent 129. Un modèle à base d'arbres ne peut pas prédire au-delà des valeurs vues à l'entraînement, ce qui le pénalise sur une série dont le niveau a fortement monté.
  \item \textbf{Ridge suit mieux la tendance} mais n'anticipe pas les chocs : son erreur atteint 59 points en 2016, et 34 points en 2020.
  \item \textbf{Les $R^2$ sont tous négatifs.} Sur le test, aucun des trois modèles ne fait mieux que la simple moyenne des valeurs du test. Le gain de Ridge sur le naïf est réel mais modeste, et il ne suffit pas à parler d'un modèle fiable.
\end{itemize}

\paragraph{Ce que l'on ne peut pas conclure.} Le test ne contient que 10 années et un seul découpage a été évalué. L'écart de 3,4 points de MAE entre Ridge et la référence naïve pourrait changer avec une autre période de test. Le classement est donc indicatif, pas définitif. De plus, une bonne performance hors échantillon n'établirait pas une relation causale.

% ======================= 6. PRÉVISION =======================
\section{Prévision exploratoire pour 2025}

Le modèle Ridge ($\alpha=0{,}1$), réentraîné sur toute la série disponible, donne pour 2025 :
\[
\hat{y}_{2025}\approx 120{,}26 \text{ points d'indice (base 100 = 2014--2016).}
\]
Cette valeur est une estimation indicative et incertaine ; elle ne représente pas une quantité en tonnes. Compte tenu des erreurs observées sur le test (MAE d'environ 22 points), il faut la lire comme un ordre de grandeur.

% ======================= 7. LIMITES =======================
\section{Limites et pistes d'amélioration}

\begin{itemize}
  \item un nombre d'années limité (64 observations, 10 pour le test) ;
  \item des chocs climatiques et commerciaux non modélisés (récoltes nationales, prix mondiaux, politiques d'importation) ;
  \item l'absence de variables explicatives externes : le modèle n'utilise que le passé de la série ;
  \item une cible exprimée en indice, qui ne renseigne pas directement sur les volumes en tonnes.
\end{itemize}
Les pistes naturelles sont l'ajout de variables comme la pluviométrie, la production nationale de blé et les prix mondiaux, ainsi qu'une évaluation sur plusieurs fenêtres temporelles successives pour stabiliser le classement des modèles.

% ======================= 8. CONCLUSION =======================
\section{Conclusion}

Ce projet répond à la question posée avec un protocole sans fuite de données : découpage chronologique, hyperparamètre choisi sur la validation, test utilisé une seule fois. Le modèle Ridge est le meilleur des trois approches comparées (MAE de 22,47 contre 25,90 pour la référence naïve), mais son avantage est modeste et les $R^2$ négatifs rappellent que la série reste difficile à prévoir à partir de son seul historique. La prévision de 120 points d'indice pour 2025 doit donc être présentée comme une estimation prudente, et non comme une certitude.

% ======================= REPRODUCTIBILITÉ =======================
\section{Reproductibilité et sources}

\begin{itemize}
  \item \textbf{Source :} FAOSTAT, domaine \emph{Trade Indices}, export du 19/11/2026 (\url{https://www.fao.org/faostat}).
  \item \textbf{Environnement :} Python, pandas, scikit-learn, matplotlib ; notebook exécuté sur Google Colab.
  \item \textbf{Graine aléatoire :} 42 pour tous les modèles.
  \item \textbf{Données simulées :} aucune, toutes les valeurs proviennent de la source.
  \item \textbf{Utilisation de l'IA :} voir le fichier \texttt{DECLARATION\_IA.md} remis avec le projet.
\end{itemize}

\end{document}
